\section{Methodology: The Kiwitsche 2-Punkt-Verfahren}
\label{sec:methodology}

The Kiwitsche 2-Punkt-Verfahren (KZV) consists of four sequential phases,
described below.
An overview is shown in \cref{fig:kzv_flowchart}.

% \begin{figure}[H]
%     \centering
%     \includegraphics[width=0.7\linewidth]{figures/kzv_flowchart.pdf}
%     \caption{Flow chart of the Kiwitsche 2-Punkt-Verfahren.}
%     \label{fig:kzv_flowchart}
% \end{figure}

\subsection{Phase 1: Bang-Bang Excitation}
\label{sec:method:phase1}

\begin{enumerate}
    \item Set the Bang-Bang controller parameters:
          switching amplitude $M = (U_{\max} - U_{\min})/2$ and hysteresis
          band $d$.
          For most processes, a hysteresis of $d \approx 0.5\text{--}2\%$ of
          the measurement range is appropriate.
    \item Set the setpoint $r$ to a representative operating point
          (e.g.\ the desired cooling temperature).
    \item Enable the Bang-Bang controller and wait until the output $y(t)$
          has settled into a \emph{stationary} limit cycle.
          Stationarity is verified by checking that three consecutive
          oscillation periods agree to within $\pm 5\%$.
\end{enumerate}

\subsection{Phase 2: Characteristic Extraction}
\label{sec:method:phase2}

Once a stationary limit cycle is established, record $N_{\min} = 5$
consecutive oscillation cycles and compute:

\begin{align}
    \Tu &= \frac{1}{N} \sum_{k=1}^{N} (t_{k+1} - t_k), \label{eq:Tu_est} \\
    \Au &= \frac{1}{N} \sum_{k=1}^{N} \bigl(y_{\max,k} - y_{\min,k}\bigr), \label{eq:Au_est}
\end{align}
where $t_k$ are the times of successive positive zero-crossings of $e(t)$,
and $y_{\max,k}$, $y_{\min,k}$ are the maximum and minimum of $y(t)$ within
each cycle.

\subsection{Phase 3: Plant Identification}
\label{sec:method:phase3}

Using the measured quantities $(\Tu, \Au, d, M)$ and the recorded time
series $(t, y, u)$, the FOPDT parameters are estimated as follows.

\paragraph{Dead-time estimate.}
For a FOPDT plant under Bang-Bang control, the output slope $\dot{y}(t)$
changes sign from negative to positive exactly at $t = t_{\mathrm{sw}} + L$
after a rising switch ($U_{\min} \to U_{\max}$), because only then does the
new input value reach the plant through the dead-time element.
The dead time is therefore estimated as
\begin{equation}
    \hat{L} = t_{\mathrm{sign-change}} - t_{\mathrm{sw}},
    \label{eq:L_est}
\end{equation}
averaged over the last $N_{\min}$ rising-switch events.

\paragraph{Time constant estimate.}
At $t = t_{\mathrm{sw,fall}} + \hat{L}$ after a falling switch ($U_{\max} \to U_{\min}$),
the FOPDT step-response slope for $U_{\min} = 0$ is
\begin{equation}
    \dot{y}\big|_{t_{\mathrm{sw,fall}} + \hat{L}}
    = \frac{K \cdot U_{\min} - y_L}{T}
    = \frac{-y_L}{T},
    \label{eq:slope_fall}
\end{equation}
where $y_L = y(t_{\mathrm{sw,fall}} + \hat{L})$.
Solving for $T$:
\begin{equation}
    \hat{T} = \frac{y_L}{|\dot{y}|_{t_{\mathrm{sw,fall}} + \hat{L}}},
    \label{eq:T_est}
\end{equation}
averaged over the last $N_{\min}$ falling-switch events.

\paragraph{Gain estimate.}
At $t = t_{\mathrm{sw,rise}} + \hat{L}$ after a rising switch, the slope is
\begin{equation}
    \dot{y}\big|_{t_{\mathrm{sw,rise}} + \hat{L}}
    = \frac{K \cdot U_{\max} - y_L}{T},
    \label{eq:slope_rise}
\end{equation}
hence
\begin{equation}
    \hat{K} = \frac{\hat{T} \cdot \dot{y}_{L,\uparrow} + y_{L,\uparrow}}{U_{\max}},
    \label{eq:K_est}
\end{equation}
averaged over the last $N_{\min}$ rising-switch events.

\begin{remark}
    \Cref{eq:T_est} is exact for $U_{\min} = 0$.
    For $U_{\min} \neq 0$, an iterative refinement involving
    $K \cdot U_{\min}$ is required; this case is not treated here.
\end{remark}

\subsection{Phase 4: PID Synthesis}
\label{sec:method:phase4}

With the identified FOPDT model $(\hat{K}, \hat{T}, \hat{L})$, apply the
IMC-PID formulas \eqref{eq:Kp}--\eqref{eq:Td} with $\lambda = \hat{L}$:

\begin{equation}
    \boxed{
    \Kp = \frac{\hat{T}}{\hat{K}(2\hat{L})},
    \quad
    \Ti = \hat{T},
    \quad
    \Td = \frac{\hat{L}}{2}.
    }
    \label{eq:pid_final}
\end{equation}

Switch the controller from Bang-Bang mode to PID mode.
No additional plant experiment is required.

\subsection{Implementation Notes}
\label{sec:method:impl}

\begin{itemize}
    \item The procedure is implemented in the Python module
          \texttt{analysis/parameter\_estimation.py}; see the repository for
          source code.
    \item All simulation and figure generation scripts are available in the
          \texttt{python/} directory; see \texttt{README.md} for usage.
    \item The same procedure can be implemented on a PLC or embedded
          controller by storing a rolling buffer of the last $3\text{--}5$
          oscillation cycles and computing the statistics online.
\end{itemize}
